% Kamp III.2, Lemma 3: the Boolean and proper-substitution closure E.

\subsubsection{Kamp's Boolean and substitution closure}
\label{sec:p09-kamp-iii2-lemma3}

Kamp next closes the current-point, basic and limiting classes under Boolean
operations and proper substitution for unary predicate letters. His
Lemma~3 states that every formula in the resulting class is expressible.
The proof invokes Lemmas~1--2 and the substitution observation R4
\cite[ch.~III, pp.~59--60, 67, 73]{Kamp1968}.
The construction below supplies that observation with its variable
conditions and gives an expression at every finite construction step.

\begin{hbdefinition}{def:p09-kamp-proper-substitution}{Proper substitution for a unary predicate}
Keep the equality-free first-order language \(L\) fixed above. Let
\(\phi\) be an \(L\)-formula, let \(P_i\) be a unary predicate letter,
and let \(\psi\) be an \(L\)-formula with exactly one free variable,
called \(v\). An \emph{alphabetic variant} changes bound variable names
and their bound occurrences without changing the binding or the free
variables. Write
\[
 \operatorname{psub}_{P_i}(\phi,\psi)
\]
for the following operation. In the fixed enumeration of \(L\)-formulas
used by Kamp, choose the first alphabetic variant \(\psi'\) of
\(\psi\) whose bound variables do not occur in \(\phi\). Replace each
original atomic occurrence \(P_i(y)\) in \(\phi\) by the result of
replacing the variable \(v\) throughout \(\psi'\) by the argument
variable \(y\), in the source's variable-replacement convention.
Leave all other predicate atoms and all order
atoms unchanged. This is Kamp's proper predicate substitution, written in
the source as \([\phi]^{\psi}_{P_i}\)
\cite[ch.~III, def.~1, p.~40]{Kamp1968}.

The enumeration here selects a formula spelling; it is a different object
from the formula class \(\operatorname{E}\) defined below. Fresh bound
names can also be chosen directly, separating any bound occurrences of
the parameter name \(v\) first. Then substitute \(y\) only for its
free occurrences, written \(\psi'[y/v]\). Any such choice gives an
alphabetic variant of the source-selected result, as proved in the annex.
Thus the construction need not identify the first entry in an unspecified
enumeration, and does not claim literal identity of selected bound names.

Substitution visits the original atoms of \(\phi\) once. It does not
then replace occurrences of \(P_i\) inside an inserted copy of
\(\psi'\). The replacement may itself contain \(P_i\), and it may
contain other unary letters already used by \(\phi\). No disjointness
or nonoccurrence condition is imposed on predicate letters. If \(P_i\)
does not occur in \(\phi\), the result is \(\phi\).

The one-free-variable condition is exact: a sentence is not an admissible
replacement here, nor is a formula with additional free parameters. The
binary relation symbol \(<\) is not a substitution target. The name
\(v\) need not be the free variable of \(\phi\); it is the argument
place of the unary replacement. When an original occurrence \(P_i(y)\)
is under a quantifier binding \(y\), that same quantifier binds the
inserted free argument occurrences of \(y\). Freshness protects the
replacement's own bound variables; it does not remove this intended binding.
\end{hbdefinition}

\begin{hbdefinition}{def:p09-kamp-temporal-substitution}{Substitution in the temporal language}
For \(\eta,\zeta\in\mathsf{TL}_1\), define
\(\eta[\zeta/q_i]\) by recursion on \(\eta\):
\[
\begin{aligned}
 q_j[\zeta/q_i]&\defeq
   \begin{cases}\zeta,&j=i,\\q_j,&j\ne i,\end{cases}\\
 (\neg\alpha)[\zeta/q_i]&\defeq\neg(\alpha[\zeta/q_i]),\\
 (\alpha\land\beta)[\zeta/q_i]
   &\defeq\alpha[\zeta/q_i]\land\beta[\zeta/q_i],\\
 S(\alpha,\beta)[\zeta/q_i]
   &\defeq S(\alpha[\zeta/q_i],\beta[\zeta/q_i]),\\
 U(\alpha,\beta)[\zeta/q_i]
   &\defeq U(\alpha[\zeta/q_i],\beta[\zeta/q_i]).
\end{aligned}
\]
The temporal language has no variable-binding constructor. At the atomic
case the output is the supplied \(\zeta\), without another substitution
inside it. The two operands of \(S,U\) retain their strict witness and
open-interval roles. Derived temporal operations are expanded into these
primitive constructors before this recursion is applied.

Disjunction is the abbreviation
\(\alpha\lor\beta\defeq\neg(\neg\alpha\land\neg\beta)\).
A finite sequence of substitutions means repeated applications of the
single operation in the stated order; the empty sequence is the identity.
It is not a simultaneous replacement or a recursive system of predicate
equations.
\end{hbdefinition}

\begin{hblemma}{thm:p09-kamp-proper-substitution}{proved=annex:p9.kamp.proper.substitution}{Proper substitution preserves expressibility}
Let \(\phi\) have exactly the free variable \(w\), and let
\(\psi\) have exactly the free variable \(v\). Suppose that
\(\eta,\zeta\in\mathsf{TL}_1\) express them over
\(\mathsf{Lin}\), respectively, with the named-variable convention of
\cref{def:p09-kamp-basic-expressibility}. Then
\(\operatorname{psub}_{P_i}(\phi,\psi)\) has exactly \(w\) free and
is expressed by \(\eta[\zeta/q_i]\). In particular,
\[
 \mathsf{Lin}\models
 \operatorname{psub}_{P_i}(\phi,\psi)
 \longleftrightarrow D(\eta[\zeta/q_i])[w/t_0].
\]
This includes \(v=w\), a predicate absent from \(\phi\), and a
replacement containing that same predicate. Each substitution is finite.
The annex also proves compatibility directly when the two first-order
inputs are the named-variable \(D\)-translations themselves. This is the
full substitution step needed from Kamp's R4
\cite[ch.~III, pp.~59--60]{Kamp1968}.
\end{hblemma}

\begin{hbdefinition}{def:p09-kamp-elementary-class}{The finite closure \(\operatorname{E}\)}
For a variable \(w\), let
\[
 \mathcal G(w)\defeq
  \operatorname{At}(w)\cup\operatorname{Bp}(w)
  \cup\operatorname{Bf}(w)\cup\operatorname{Lp}(w)
  \cup\operatorname{Lf}(w).
\]
Here \(\operatorname{At}\) is the unstarred source class from
\cref{def:p09-kamp-iii1-components}; it consists of finite nonempty
conjunctions of unary literals and the particular tautology
\(P_0(w)\lor\neg P_0(w)\). The basic and limiting classes are those
of \cref{def:p09-kamp-basic-classes,def:p09-kamp-limiting-classes}.

Kamp defines \(\operatorname{E}(w)\) as the closure of this union under
negation, conjunction, disjunction and proper predicate substitution
\cite[ch.~III, p.~67]{Kamp1968}. Explicit finite stages for that least
closure are as follows. Set \(\operatorname{E}_0(w)\defeq\mathcal G(w)\).
For each \(m\in\mathbb N\), at stage \(m+1\) retain
\(\operatorname{E}_m(w)\) and add:
\[
\begin{gathered}
 \neg\phi,\qquad \phi\land\psi,\qquad \phi\lor\psi,
 \qquad \operatorname{psub}_{P_i}(\phi,\psi),\\
 \phi,\psi\in\operatorname{E}_m(w),\qquad i\in\mathbb N.
\end{gathered}
\]
Put
\[
 \operatorname{E}(w)\defeq\bigcup_{m\geq0}\operatorname{E}_m(w),
 \qquad
 \operatorname{E}\defeq\bigcup_{j\geq0}\operatorname{E}(t_j).
\]
The substitution operands in this fixed-variable closure both belong to
\(\operatorname{E}_m(w)\); the preceding substitution lemma permits the
more general parameter name \(v\) as well. Each constructed formula has
exactly its named variable free. The stages describe finite constructions,
not a closure under semantic equivalence, arbitrary quantification,
infinitary Boolean operations, or substitution of arbitrary formulas from
outside the class. A formula can have several construction histories.

There are no primitive zero-operand Boolean constructors in this grammar.
When an empty finite conjunction or disjunction is requested as an
abbreviation, use respectively
\(T^K_w=P_0(w)\lor\neg P_0(w)\) or
\(F^K_w=P_0(w)\land\neg P_0(w)\), which already belong to
\(\operatorname{At}(w)\). This uses the declared source letter \(P_0\).
\end{hbdefinition}

\begin{hbmethod}{def:p09-kamp-elementary-translation}{A temporal expression from a finite construction}
The input is a finite construction of \(\phi\in\operatorname{E}(w)\),
including the constituent basic or limiting pattern at a seed. Read that
construction from its leaves upward and assign a temporal formula to each
node:
\begin{enumerate}
\item For an \(\operatorname{At}(w)\) seed, replace \(P_i(w)\) by
      \(q_i\) and preserve its literal negations and conjunctions.
      For its particular tautology use
      \(\neg(q_0\land\neg q_0)\).
\item For a basic past or future seed, use respectively the \(L_n\) or
      \(R_0\) construction in \cref{thm:p09-kamp-iii2-lemma1}.
      For a limiting past or future seed, use respectively
      \(\Lambda^-_\sigma\) or \(\Lambda^+_\sigma\) from
      \cref{thm:p09-kamp-iii2-lemma2}, including their one-letter seeds.
\item For a negation node with child expression \(\eta\), return
      \(\neg\eta\). For conjunction and disjunction nodes with child
      expressions \(\eta,\zeta\), return respectively
      \(\eta\land\zeta\) and
      \(\neg(\neg\eta\land\neg\zeta)\).
\item For a node \(\operatorname{psub}_{P_i}(\phi_1,\phi_2)\),
      first obtain its two child expressions \(\eta,\zeta\), then
      return \(\eta[\zeta/q_i]\) by the finite structural recursion
      just defined.
\end{enumerate}
Stage retention introduces no new operation. The output is one finite
\(\mathsf{TL}_1\)-formula. A construction history is an explicit witness
of membership; this method does not assert a decision procedure for
membership of an arbitrary first-order formula in \(\operatorname{E}\).
\end{hbmethod}

\begin{hblemma}{thm:p09-kamp-iii2-lemma3}{proved=annex:p9.kamp.iii2.lemma3}{Kamp's expressibility lemma for \(\operatorname{E}\)}
Every member of \(\operatorname{E}\) is expressible over
\(\mathsf{Lin}\). More explicitly, if a finite construction of
\(\phi\in\operatorname{E}(w)\) returns \(\eta\) by
\cref{def:p09-kamp-elementary-translation}, then
\[
             \mathsf{Lin}\models
                  \phi\longleftrightarrow D(\eta)[w/t_0].
\]
The conclusion holds for every interpretation of the unary predicates
on every nonempty strict linear order. The output uses only proposition
letters corresponding to unary predicate letters of the constructed
formula, including the source's \(P_0\) when a source truth/falsehood
seed contains it. No order-completeness, density, or endpointlessness
assumption is used \cite[ch.~III, lemma~3, p.~73]{Kamp1968}.
\end{hblemma}

\begin{hbexample}{def:p09-kamp-substitution-example}{A limiting predicate inside a basic future pattern}
Let \(A,B,C\) be distinct unary predicate letters. Take the basic future
formula with word \((A,B,C)\), expressed by
\(\widehat A\land U(\widehat C,\widehat B)\), and substitute
\(\partial^-C(w)\) properly for \(C\). Both operands belong to
\(\operatorname{E}(w)\). The resulting temporal formula is
\[
                  \widehat A\land U(P'(\widehat C),\widehat B).
\]
The inserted limit is tested at the future endpoint chosen by \(U\),
not at the original current point. Its internal occurrence of \(C\)
is retained: a single substitution step does not call itself again inside
the inserted formula. The annex writes the first-order result with fresh
variables and verifies this example.
\end{hbexample}

\paragraph{The following order-completeness boundary.}
Immediately after Lemma~3, Kamp changes the default interpretation of
``valid'', ``equivalent'' and ``expressible'' from \(\mathsf{Lin}\) to
\(\mathsf{Com}\). The next result is Lemma~4, concerning conjunction
for the interval property \(I(x,y)\); its proof begins on printed p.~73.
Its elementary interval classes, interval property and remaining cases
are the next separate dependency. The present closure lemma supplies no
proof that every one-variable first-order formula belongs to
\(\operatorname{E}\), and does not complete Theorem~III.2
\cite[ch.~III, pp.~67--68, 73]{Kamp1968}.
